% Options for packages loaded elsewhere
\PassOptionsToPackage{unicode}{hyperref}
\PassOptionsToPackage{hyphens}{url}
\PassOptionsToPackage{dvipsnames,svgnames,x11names}{xcolor}
%
\documentclass[
  letterpaper,
]{article}

\usepackage{amsmath,amssymb}
\usepackage{iftex}
\ifPDFTeX
  \usepackage[T1]{fontenc}
  \usepackage[utf8]{inputenc}
  \usepackage{textcomp} % provide euro and other symbols
\else % if luatex or xetex
  \usepackage{unicode-math}
  \defaultfontfeatures{Scale=MatchLowercase}
  \defaultfontfeatures[\rmfamily]{Ligatures=TeX,Scale=1}
\fi
\usepackage{lmodern}
\ifPDFTeX\else  
    % xetex/luatex font selection
\fi
% Use upquote if available, for straight quotes in verbatim environments
\IfFileExists{upquote.sty}{\usepackage{upquote}}{}
\IfFileExists{microtype.sty}{% use microtype if available
  \usepackage[]{microtype}
  \UseMicrotypeSet[protrusion]{basicmath} % disable protrusion for tt fonts
}{}
\makeatletter
\@ifundefined{KOMAClassName}{% if non-KOMA class
  \IfFileExists{parskip.sty}{%
    \usepackage{parskip}
  }{% else
    \setlength{\parindent}{0pt}
    \setlength{\parskip}{6pt plus 2pt minus 1pt}}
}{% if KOMA class
  \KOMAoptions{parskip=half}}
\makeatother
\usepackage{xcolor}
\usepackage[margin=1in]{geometry}
\setlength{\emergencystretch}{3em} % prevent overfull lines
\setcounter{secnumdepth}{5}
% Make \paragraph and \subparagraph free-standing
\makeatletter
\ifx\paragraph\undefined\else
  \let\oldparagraph\paragraph
  \renewcommand{\paragraph}{
    \@ifstar
      \xxxParagraphStar
      \xxxParagraphNoStar
  }
  \newcommand{\xxxParagraphStar}[1]{\oldparagraph*{#1}\mbox{}}
  \newcommand{\xxxParagraphNoStar}[1]{\oldparagraph{#1}\mbox{}}
\fi
\ifx\subparagraph\undefined\else
  \let\oldsubparagraph\subparagraph
  \renewcommand{\subparagraph}{
    \@ifstar
      \xxxSubParagraphStar
      \xxxSubParagraphNoStar
  }
  \newcommand{\xxxSubParagraphStar}[1]{\oldsubparagraph*{#1}\mbox{}}
  \newcommand{\xxxSubParagraphNoStar}[1]{\oldsubparagraph{#1}\mbox{}}
\fi
\makeatother


\providecommand{\tightlist}{%
  \setlength{\itemsep}{0pt}\setlength{\parskip}{0pt}}\usepackage{longtable,booktabs,array}
\usepackage{calc} % for calculating minipage widths
% Correct order of tables after \paragraph or \subparagraph
\usepackage{etoolbox}
\makeatletter
\patchcmd\longtable{\par}{\if@noskipsec\mbox{}\fi\par}{}{}
\makeatother
% Allow footnotes in longtable head/foot
\IfFileExists{footnotehyper.sty}{\usepackage{footnotehyper}}{\usepackage{footnote}}
\makesavenoteenv{longtable}
\usepackage{graphicx}
\makeatletter
\newsavebox\pandoc@box
\newcommand*\pandocbounded[1]{% scales image to fit in text height/width
  \sbox\pandoc@box{#1}%
  \Gscale@div\@tempa{\textheight}{\dimexpr\ht\pandoc@box+\dp\pandoc@box\relax}%
  \Gscale@div\@tempb{\linewidth}{\wd\pandoc@box}%
  \ifdim\@tempb\p@<\@tempa\p@\let\@tempa\@tempb\fi% select the smaller of both
  \ifdim\@tempa\p@<\p@\scalebox{\@tempa}{\usebox\pandoc@box}%
  \else\usebox{\pandoc@box}%
  \fi%
}
% Set default figure placement to htbp
\def\fps@figure{htbp}
\makeatother
% definitions for citeproc citations
\NewDocumentCommand\citeproctext{}{}
\NewDocumentCommand\citeproc{mm}{%
  \begingroup\def\citeproctext{#2}\cite{#1}\endgroup}
\makeatletter
 % allow citations to break across lines
 \let\@cite@ofmt\@firstofone
 % avoid brackets around text for \cite:
 \def\@biblabel#1{}
 \def\@cite#1#2{{#1\if@tempswa , #2\fi}}
\makeatother
\newlength{\cslhangindent}
\setlength{\cslhangindent}{1.5em}
\newlength{\csllabelwidth}
\setlength{\csllabelwidth}{3em}
\newenvironment{CSLReferences}[2] % #1 hanging-indent, #2 entry-spacing
 {\begin{list}{}{%
  \setlength{\itemindent}{0pt}
  \setlength{\leftmargin}{0pt}
  \setlength{\parsep}{0pt}
  % turn on hanging indent if param 1 is 1
  \ifodd #1
   \setlength{\leftmargin}{\cslhangindent}
   \setlength{\itemindent}{-1\cslhangindent}
  \fi
  % set entry spacing
  \setlength{\itemsep}{#2\baselineskip}}}
 {\end{list}}
\usepackage{calc}
\newcommand{\CSLBlock}[1]{\hfill\break\parbox[t]{\linewidth}{\strut\ignorespaces#1\strut}}
\newcommand{\CSLLeftMargin}[1]{\parbox[t]{\csllabelwidth}{\strut#1\strut}}
\newcommand{\CSLRightInline}[1]{\parbox[t]{\linewidth - \csllabelwidth}{\strut#1\strut}}
\newcommand{\CSLIndent}[1]{\hspace{\cslhangindent}#1}

\makeatletter
\@ifpackageloaded{tcolorbox}{}{\usepackage[skins,breakable]{tcolorbox}}
\@ifpackageloaded{fontawesome5}{}{\usepackage{fontawesome5}}
\definecolor{quarto-callout-color}{HTML}{909090}
\definecolor{quarto-callout-note-color}{HTML}{0758E5}
\definecolor{quarto-callout-important-color}{HTML}{CC1914}
\definecolor{quarto-callout-warning-color}{HTML}{EB9113}
\definecolor{quarto-callout-tip-color}{HTML}{00A047}
\definecolor{quarto-callout-caution-color}{HTML}{FC5300}
\definecolor{quarto-callout-color-frame}{HTML}{acacac}
\definecolor{quarto-callout-note-color-frame}{HTML}{4582ec}
\definecolor{quarto-callout-important-color-frame}{HTML}{d9534f}
\definecolor{quarto-callout-warning-color-frame}{HTML}{f0ad4e}
\definecolor{quarto-callout-tip-color-frame}{HTML}{02b875}
\definecolor{quarto-callout-caution-color-frame}{HTML}{fd7e14}
\makeatother
\makeatletter
\@ifpackageloaded{caption}{}{\usepackage{caption}}
\AtBeginDocument{%
\ifdefined\contentsname
  \renewcommand*\contentsname{Table of contents}
\else
  \newcommand\contentsname{Table of contents}
\fi
\ifdefined\listfigurename
  \renewcommand*\listfigurename{List of Figures}
\else
  \newcommand\listfigurename{List of Figures}
\fi
\ifdefined\listtablename
  \renewcommand*\listtablename{List of Tables}
\else
  \newcommand\listtablename{List of Tables}
\fi
\ifdefined\figurename
  \renewcommand*\figurename{Figure}
\else
  \newcommand\figurename{Figure}
\fi
\ifdefined\tablename
  \renewcommand*\tablename{Table}
\else
  \newcommand\tablename{Table}
\fi
}
\@ifpackageloaded{float}{}{\usepackage{float}}
\floatstyle{ruled}
\@ifundefined{c@chapter}{\newfloat{codelisting}{h}{lop}}{\newfloat{codelisting}{h}{lop}[chapter]}
\floatname{codelisting}{Listing}
\newcommand*\listoflistings{\listof{codelisting}{List of Listings}}
\makeatother
\makeatletter
\makeatother
\makeatletter
\@ifpackageloaded{caption}{}{\usepackage{caption}}
\@ifpackageloaded{subcaption}{}{\usepackage{subcaption}}
\makeatother

\usepackage{bookmark}

\IfFileExists{xurl.sty}{\usepackage{xurl}}{} % add URL line breaks if available
\urlstyle{same} % disable monospaced font for URLs
\hypersetup{
  pdftitle={Proximal-to-Distal Energy Transfer in the Golf Swing: Where It Occurs, How It Differs with Skill, and When It Should Happen},
  pdfauthor={Dieter Olson},
  pdfkeywords={golf swing, proximal-to-distal sequencing, kinematic
sequence, energy transfer, counterfactual dynamics, ZTCF, ZVCF},
  colorlinks=true,
  linkcolor={blue},
  filecolor={Maroon},
  citecolor={Blue},
  urlcolor={Blue},
  pdfcreator={LaTeX via pandoc}}


\title{Proximal-to-Distal Energy Transfer in the Golf Swing: Where It
Occurs, How It Differs with Skill, and When It Should Happen}
\usepackage{etoolbox}
\makeatletter
\providecommand{\subtitle}[1]{% add subtitle to \maketitle
  \apptocmd{\@title}{\par {\large #1 \par}}{}{}
}
\makeatother
\subtitle{A literature synthesis and modeling perspective grounded in
the UpstreamDrift ZTCF/ZVCF counterfactual framework}
\author{Dieter Olson}
\date{2026-08-07}

\begin{document}
\maketitle
\begin{abstract}
Proximal-to-distal (P→D) sequencing --- the ordered peaking of pelvis,
thorax, arm, and club rotational speeds --- is the most widely reported
organizing principle of the golf downswing. It is usually interpreted as
evidence that mechanical energy flows outward along the kinetic chain,
and it is often translated into the coaching prescription that energy
should be passed distally as early and as completely as possible. This
article reviews the mechanics and the empirical literature behind that
interpretation and asks a sharper question: \emph{when} during the
downswing should the P→D transfer occur? Double-pendulum analyses,
delayed-release simulations, and optimal-control results consistently
favor accumulating energy proximally first (the arms-plus-club system
rotating as a quasi-rigid body), with the transfer to the club
deliberately \emph{restrained} early in the downswing and
\emph{accelerated} late. Skill-level comparisons show that experts
differ from amateurs more in transition timing, segment-speed
magnitudes, and transfer efficiency than in the mere presence of a P→D
order, and methodological work cautions that peak-ordering alone is a
weak proxy for energy flow. We connect this literature to the
zero-torque counterfactual (ZTCF) and zero-velocity counterfactual
(ZVCF) decompositions implemented in the UpstreamDrift platform, which
separate, at every instant of a measured swing, the acceleration
attributable to accumulated motion (the drift term) from the
acceleration attributable to the applied joint torques (the control
term). We formalize the ``retain-early / release-late'' hypothesis in
counterfactual terms, inventory the repository's models and
energy-transfer computations, identify the measurement gaps (notably
segmental energy-flow accounting), and specify falsifiable simulation
experiments --- torque-timing sweeps, counterfactual energy budgets,
re-simulation parity --- that can decide the question quantitatively. No
new experimental data are reported; every quantitative claim is sourced
to the cited literature or to inspectable code.
\end{abstract}

\renewcommand*\contentsname{Table of contents}
{
\hypersetup{linkcolor=}
\setcounter{tocdepth}{3}
\tableofcontents
}

\begin{tcolorbox}[enhanced jigsaw, arc=.35mm, toprule=.15mm, breakable, toptitle=1mm, opacityback=0, coltitle=black, bottomrule=.15mm, left=2mm, title=\textcolor{quarto-callout-note-color}{\faInfo}\hspace{0.5em}{Provenance and scope}, titlerule=0mm, colbacktitle=quarto-callout-note-color!10!white, colback=white, bottomtitle=1mm, colframe=quarto-callout-note-color-frame, leftrule=.75mm, opacitybacktitle=0.6, rightrule=.15mm]

Deliverable of UpstreamDrift epic
\href{https://github.com/D-sorganization/UpstreamDrift/issues/8373}{\#8373}
(sub-issues \#8374--\#8379). A \emph{review and modeling-perspective}
paper: no new experimental data are reported; the experiments of
Section~\ref{sec-proposed-experiments} are planned work tracked in issue
\href{https://github.com/D-sorganization/UpstreamDrift/issues/8377}{\#8377}.
All references were verified against publication records at the time of
writing; repository links point to the \texttt{main} branch of
\href{https://github.com/D-sorganization/UpstreamDrift}{D-sorganization/UpstreamDrift}.

\end{tcolorbox}

\section{Introduction}\label{sec-intro}

Clubhead speed at impact is the dominant controllable determinant of
drive distance, and the golf downswing is one of the most-studied
examples of a \emph{sequenced} multi-segment rotation in sport
biomechanics (\citeproc{ref-hume2005}{Hume, Keogh, and Reid 2005};
\citeproc{ref-mcphee2022}{McPhee 2022}). Across measurement systems and
populations, the downswing of a skilled golfer shows a characteristic
order: the pelvis reaches its peak angular speed first, then the thorax,
then the lead arm, and finally the club, with each successive segment
reaching a higher peak speed than the one before it
(\citeproc{ref-cheetham2008}{Cheetham et al. 2008};
\citeproc{ref-tinmark2010}{Tinmark et al. 2010}). This pattern --- the
\emph{kinematic sequence}, an instance of the general proximal-to-distal
(P→D) sequencing observed in throwing and striking skills
(\citeproc{ref-putnam1993}{Putnam 1993}) --- is routinely interpreted as
the visible signature of mechanical energy being generated in the large
proximal segments and transferred outward to the small, fast distal
ones.

That interpretation immediately raises a question that the descriptive
literature does not settle:

\begin{quote}
\textbf{Should a golfer attempt to transfer energy from proximal to
distal segments \emph{early} in the downswing, or should energy be
\emph{retained} in the proximal segments early and the transfer
deliberately \emph{accelerated} late?}
\end{quote}

The question is not academic. The two strategies imply opposite coaching
cues (``fire the club from the top'' versus ``hold the lag and let it
release''), opposite torque programs, and different predictions for what
counterfactual decompositions of a measured swing should show. It also
connects to a long-standing puzzle: double-pendulum analyses have argued
since the 1960s that the highest clubhead speeds arise when the distal
joint is \emph{not} driven early --- when release is passively or
actively delayed (\citeproc{ref-cochran1968}{Cochran and Stobbs 1968};
\citeproc{ref-jorgensen1970}{T. Jorgensen 1970};
\citeproc{ref-milburn1982}{Milburn 1982};
\citeproc{ref-pickering1999}{Pickering and Vickers 1999}) --- while the
coaching translation of the kinematic sequence often emphasizes
initiating distal motion promptly after transition.

The UpstreamDrift project approaches this question with
\emph{counterfactual dynamics}. Two decompositions are implemented
across its physics engines:

\begin{itemize}
\tightlist
\item
  the \textbf{zero-torque counterfactual (ZTCF)} --- the acceleration
  the system would exhibit at the current instant if all applied joint
  torques were removed, isolating what accumulated motion, gravity, and
  inter-segment coupling produce on their own; and
\item
  the \textbf{zero-velocity counterfactual (ZVCF)} --- the acceleration
  the applied torques, together with gravity, would produce at the
  current pose if the system were momentarily at rest, removing all
  motion-dependent effects (the pure torque contribution is the control
  term of the drift-control split defined in Section~\ref{sec-ztcf}).
\end{itemize}

Preliminary ZTCF/ZVCF analyses with the repository's MATLAB pipeline
(Section~\ref{sec-ztcf}) motivated the hypothesis examined here: that
\emph{restraining} the P→D handoff early in the downswing --- keeping
kinetic energy in the proximal segments and in the arms-plus-club system
rotating as a quasi-rigid unit --- and accelerating the handoff late may
be the more effective strategy. This article evaluates that hypothesis
against the published literature and specifies how to test it
quantitatively with the tools already in the repository.

\textbf{Contributions.}

\begin{enumerate}
\def\labelenumi{\arabic{enumi}.}
\tightlist
\item
  a synthesis of the mechanics of sequenced swings, from the
  summation-of-speed principle to segment-interaction dynamics and
  parametric energy transfer (Section~\ref{sec-mechanics});
\item
  a methodological account of what ``energy transfer'' means
  operationally --- joint power, joint-force power, segmental energy
  accounting, and induced accelerations (Section~\ref{sec-energyflow});
\item
  a critical review of the empirical evidence on where P→D transfer
  occurs and how it differs between skilled and less-skilled players
  (Section~\ref{sec-evidence});
\item
  a formalization of the early-versus-late transfer question, including
  the ZTCF/ZVCF counterfactual framework and the specific predictions
  the ``retain-early / release-late'' hypothesis makes
  (Section~\ref{sec-timing}, Section~\ref{sec-ztcf}); and
\item
  an inventory of the models and computations available in UpstreamDrift
  today, an honest gap analysis, and a set of falsifiable simulation
  experiments that would decide the question (Section~\ref{sec-repo},
  Section~\ref{sec-proposed-experiments}).
\end{enumerate}

\section{Mechanics of sequenced swings}\label{sec-mechanics}

\subsection{The summation-of-speed principle and its
limits}\label{sec-summation}

The oldest formalization of P→D sequencing is the \emph{summation of
speed} principle: in skills demanding high end-point speed, each segment
should begin its motion at the moment of peak speed of the segment
proximal to it, so that distal speed is ``built on top of'' proximal
speed (\citeproc{ref-bunn1972}{Bunn 1972}). Milburn applied the
principle directly to golf with a double-pendulum analysis of the
downswing and found the characteristic pattern of an initially dominant
proximal (arm) rotation followed by rapid distal (club) rotation, noting
that \emph{``an initial delay in the uncocking of the wrist''} allows
the arm to reach greater acceleration and the club's acceleration to be
summed with the proximal segment's maximum
(\citeproc{ref-milburn1982}{Milburn 1982}). Note the structure of that
finding: the summation principle, taken naively, prescribes
\emph{sequential initiation}; Milburn's analysis showed the benefit
comes specifically from \emph{delaying} the distal contribution.

Putnam's segment-interaction analyses put the principle on a dynamical
footing and exposed its limits (\citeproc{ref-putnam1991}{Putnam 1991},
\citeproc{ref-putnam1993}{1993}). In a linked chain, each segment's
angular acceleration depends not only on joint torques but on
\emph{motion-dependent interaction torques} --- centripetal and Coriolis
terms proportional to squares and products of segment angular
velocities. Putnam showed the P→D pattern in striking and throwing
arises substantially from these interaction torques: proximal
deceleration and distal acceleration late in the motion are coupled
through the joint force and occur even without a distal joint torque.
Two implications matter for golf. First, kinematic order does not reveal
\emph{which} torques produced it: an identical peak sequence can arise
from different mixes of muscular and interaction torques. Second, the
proximal deceleration observed in experts need not be an active
``braking'' strategy --- it can be the dynamical \emph{consequence} of
the distal segment accelerating and back-reacting on the chain. Marshall
and Elliott further cautioned that planar peak-ordering analyses miss
long-axis rotations, which peak very late in racquet skills and
contribute substantially to end-point speed
(\citeproc{ref-marshall2000}{Marshall and Elliott 2000}); the golf
analogue is forearm rotation and wrist action in the final phase of the
downswing.

\subsection{Double-pendulum mechanics of the golf
downswing}\label{sec-doublependulum}

The two-lever, one-hinge model of the downswing --- an ``arm'' link
rotating about a fixed or moving hub, connected by a wrist hinge to a
``club'' link --- has been the workhorse of golf mechanics since Cochran
and Stobbs' \emph{Search for the Perfect Swing}
(\citeproc{ref-cochran1968}{Cochran and Stobbs 1968}) and Jorgensen's
dynamical analysis fitted to stroboscopic photographs of a professional
swing (\citeproc{ref-jorgensen1970}{T. Jorgensen 1970};
\citeproc{ref-jorgensen1999}{T. P. Jorgensen 1999}). Its enduring value
here is that it makes the \emph{timing} of the P→D handoff an explicit
variable (the ``release'': the point at which the wrist angle is
allowed, or driven, to open) and exposes the mechanism by which delay
pays. Three results anchor the analysis:

\textbf{Delayed release increases clubhead speed.} Pickering and
Vickers, sweeping release angles in a driven double-pendulum model,
found clubhead velocity at impact increases as release is delayed ---
the classical ``late hit'' (\citeproc{ref-pickering1999}{Pickering and
Vickers 1999}). Jorgensen's fitted model similarly reproduced the expert
swing only with a substantially delayed opening of the wrist angle
(\citeproc{ref-jorgensen1970}{T. Jorgensen 1970};
\citeproc{ref-jorgensen1999}{T. P. Jorgensen 1999}). In these models an
\emph{early} release (early P→D handoff) is strictly inferior: club
kinetic energy purchased early is paid for with arm speed that never
develops.

\textbf{Late wrist torque helps; early wrist torque is the wrong tool.}
Sprigings and Neal, using a torque-driven three-segment model with
physiologically bounded actuators, found that an active wrist torque
applied \emph{during the latter stages of the downswing} increased
clubhead speed by approximately 4\%
(\citeproc{ref-sprigings2000}{Sprigings and Neal 2000}); Sprigings and
MacKenzie confirmed the delayed-release benefit in a follow-up
simulation focused specifically on the timing of release
(\citeproc{ref-sprigings2002}{Sprigings and MacKenzie 2002}). The point
is not that distal effort is useless, but that its productive window is
late.

\textbf{Energy transfer to the club can be largely passive --- and is
governed by the changing geometry.} White analyzed an \emph{undriven}
double pendulum and showed that the transfer of kinetic energy from arms
to club late in the downswing is a parametric-transfer phenomenon: as
the club swings out, the effective moment of inertia of the system
changes, and conservation of energy and angular momentum move kinetic
energy distally without any wrist torque at all
(\citeproc{ref-white2006}{White 2006}). The efficiency of this transfer
is controlled principally by the wrist-cock angle maintained before
release (\citeproc{ref-white2006}{White 2006}). Miura's analysis of the
``inward pull'' complements this: an expert's hands moving toward the
body near impact --- equivalent to shrinking the hub radius when
centrifugal loading is fully developed --- performs work against the
centrifugal force that appears as additional clubhead kinetic energy
(\citeproc{ref-miura2001}{Miura 2001}). Nesbit and McGinnis' hub-path
analyses generalize the same idea: the non-circular hand path of skilled
golfers, including its late-downswing tightening, is a mechanism for
shaping kinetic transfer from golfer to club
(\citeproc{ref-nesbit2009hub}{Nesbit and McGinnis 2009};
\citeproc{ref-nesbit2014hub}{Nesbit and McGinnis 2014}).

Together these results describe a machine whose highest-output mode is:
\emph{load the proximal rotation first, hold the distal joint folded
(retaining energy in a compact, fast-rotating configuration), then let
geometry and --- in the last phase --- modest distal torque move the
energy outward rapidly}. The double-pendulum literature is already an
argument for \textbf{retain-early / release-late}.

\subsection{Optimal control and torque
sequencing}\label{sec-optimalcontrol}

If the timing question is posed as an optimization --- what torque
program maximizes clubhead speed (or carry) subject to actuator limits?
--- the answer has been studied repeatedly. Lampsa's early
optimal-control study of a driven double pendulum computed torque
histories maximizing drive distance, and notably found the
\emph{optimal} torques differed in structure from those measured (via
inverse dynamics) in his own swing (\citeproc{ref-lampsa1975}{Lampsa
1975}) --- early evidence that human swings need not be optimal and that
the optimum is a meaningful target rather than a redescription of expert
behavior. Sharp, fitting shoulder-arm-club models to expert swing data
and then optimizing torque programs, established a pattern for the best
use of the available driving torques across the downswing phases, with
the distal contribution structured late rather than distributed evenly
(\citeproc{ref-sharp2009}{Sharp 2009}). Three-dimensional
forward-dynamics models extended the approach with anatomically
meaningful actuators and flexible shafts: MacKenzie and Sprigings' 3D
model (\citeproc{ref-mackenzie2009}{MacKenzie and Sprigings 2009}), the
carry-optimized model of Balzerson, Banerjee and McPhee
(\citeproc{ref-balzerson2016}{Balzerson, Banerjee, and McPhee 2016}),
and the six-degree-of-freedom dynamic optimization of McNally and McPhee
(\citeproc{ref-mcnally2018}{McNally and McPhee 2018}) all generate
expert-like clubhead speeds with torque programs whose distal elements
act late in the downswing. McPhee's review surveys this modeling lineage
(\citeproc{ref-mcphee2022}{McPhee 2022}).

A caution cuts the other way: optimal control optimizes \emph{within a
model}, and the sensitivity of optimal timing to model structure (hub
mobility, shaft flexibility, actuator rate limits) is exactly why a
multi-engine platform with cross-engine validation
(Section~\ref{sec-repo}) is the right instrument for revisiting the
question.

\section{What ``energy transfer'' means
operationally}\label{sec-energyflow}

Claims about energy transfer require an accounting framework. Four
levels of analysis appear in the literature, and they are not
interchangeable.

\subsection{Joint power and work}\label{joint-power-and-work}

For a joint spanning segments \(i\) and \(i{+}1\) with net joint torque
\(\tau\) and relative angular velocity \(\omega\), the \emph{joint
(moment) power} is \(P_m = \tau\,\omega\); its time integral over a
phase is the mechanical work done by the net moment. Positive work
indicates generation (or concentric action), negative work absorption.
This is the quantity most golf studies report, and it is the basis of
the work-power budget of Nesbit and Serrano
(Section~\ref{sec-workbudget}) (\citeproc{ref-nesbit2005work}{Nesbit and
Serrano 2005}).

\subsection{Joint-force power and segmental energy
accounting}\label{sec-rw}

Joint moments are not the only conduit of mechanical energy: the joint
reaction force \(\mathbf{F}\) acting at a joint moving with velocity
\(\mathbf{v}\) transmits \emph{joint-force power}
\(P_f = \mathbf{F}\cdot\mathbf{v}\) from one segment to the next.
Robertson and Winter's classic formulation balances, for each segment,
the rate of change of segment mechanical energy against the sum of
joint-force powers and moment powers at its proximal and distal ends,
distinguishing energy \emph{transferred between} segments from energy
\emph{generated or absorbed} at joints
(\citeproc{ref-robertson1980}{Robertson and Winter 1980}). This is the
correct level of analysis for the question posed here, because the
``retain-early / release-late'' hypothesis is precisely a claim about
the \emph{time course of inter-segment transfer terms}, not about peak
angular velocities. As Section~\ref{sec-repo-gaps} notes, this
accounting is the principal measurement gap in the UpstreamDrift Python
stack today (the MATLAB pipeline's per-joint linear-plus-angular work
decomposition is its closest existing realization).

\subsection{Induced accelerations and dynamic
coupling}\label{induced-accelerations-and-dynamic-coupling}

In a coupled chain, any single torque induces accelerations at
\emph{every} joint --- including joints it does not span --- through the
inverse inertia matrix (\citeproc{ref-zajac1989}{Zajac and Gordon
1989}). Induced-acceleration and power-transfer analyses in locomotion
research (\citeproc{ref-zajac2002}{Zajac, Neptune, and Kautz 2002}) show
that intuitions of the form ``muscle X accelerates joint Y'' routinely
fail in coupled systems. For golf this matters because a torso torque
early in the downswing partly accelerates the \emph{club} (through
coupling), and a wrist torque late in the downswing back-reacts on the
torso; the net energy routing is an emergent property of the whole
chain. The ZTCF/ZVCF decompositions of Section~\ref{sec-ztcf} are the
counterfactual cousins of this analysis: rather than attributing
acceleration to individual torques, they split the total acceleration
into a motion-dependent (drift) part and a control part, instant by
instant.

\subsection{Kinematic sequence as a proxy --- and its failure
modes}\label{kinematic-sequence-as-a-proxy-and-its-failure-modes}

Peak-ordering of segment angular speeds is a \emph{proxy} for energy
flow, not a measurement of it. Marsan and colleagues showed that the
identified kinematic sequence in golf depends materially on which
angular-velocity component (or norm, or projection) is analyzed:
thirteen golfers analyzed with seven component-selection methods yielded
different sequences (\citeproc{ref-marsan2019}{Marsan et al. 2019}).
Marshall and Elliott's long-axis-rotation caution
(\citeproc{ref-marshall2000}{Marshall and Elliott 2000}) is a special
case of the same problem. Any serious test of the timing hypothesis must
therefore be run at the level of energy and power accounting
(Section~\ref{sec-rw}), with the kinematic sequence reported as
corroborating description only.

\section{Empirical evidence in golfers}\label{sec-evidence}

\subsection{The work and power budget of the
swing}\label{sec-workbudget}

Nesbit's full-body kinematic-kinetic study and the companion work-power
analysis of Nesbit and Serrano remain the most complete published energy
budgets of the swing (\citeproc{ref-nesbit2005kinematic}{Nesbit 2005};
\citeproc{ref-nesbit2005work}{Nesbit and Serrano 2005}). Using
subject-specific full-body models of four golfers of widely differing
skill, they report that the trunk complex (lumbar and thoracic spine
plus hips) contributes 68.7--72.2\% of total swing work, the shoulder
and arm joints 24.3--28\%, and the legs only 3.3--3.8\%
(\citeproc{ref-nesbit2005work}{Nesbit and Serrano 2005}). Two further
observations from the same study bear directly on this article's
question. First, power generation and clubhead speed decreased
systematically with increasing handicap. Second, the authors' phase
analysis locates the \emph{generation} of energy predominantly in the
earlier downswing (torso-dominated) and the \emph{conversion} into
clubhead speed late (shoulder/arm and club interface) --- i.e.,
generation and transfer are temporally separated, which is a
precondition for the retain-early / release-late strategy to even be
possible.

The ground is the ultimate source of the external impulse. Han and
colleagues characterized golfer-ground interaction in 63 skilled golfers
and found meaningful associations between ground-interaction force and
moment parameters and maximum clubhead speed (\citeproc{ref-han2019}{Han
et al. 2019}). Recent work has pushed this to an explicitly energetic
formulation: in thirty golfers (fifteen professionals, fifteen
high-level amateurs), foot-ground interaction variables alone explained
35.5\% of clubhead-speed variance, rising to 75.4\% when trunk
sequencing and an \emph{impulse-based energy-transfer efficiency}
measure (effectiveness of mechanical impulse transmission from proximal
segments to the club) were added; mediation analysis attributed the
dominant indirect pathway to transfer efficiency rather than to
sequencing per se
(\citeproc{ref-frontiers2026footground}{{``Foot--Ground Interaction and
Clubhead Speed: Impulse-Based Energy Transfer as the Key Mechanism in
the Golf Swing''} 2026}). That last distinction --- transfer
\emph{efficiency} mattering where sequence \emph{order} alone does not
--- is a recurring theme of the skill-level literature reviewed next.

\subsection{Kinematic-sequence evidence and skill-level
differences}\label{sec-skill}

The descriptive P→D findings are robust in skilled populations. Tinmark
and colleagues, studying elite golfers across full and partial shots,
found a significant proximal-to-distal temporal ordering of peak speeds
(pelvis → torso → hands/club) with successive amplification of peak
speed at every level of effort, in both men and women
(\citeproc{ref-tinmark2010}{Tinmark et al. 2010}). Cheetham and
colleagues compared amateur recreational players with touring
professionals on transition- and downswing-phase kinematic sequence
parameters and found professionals showed higher peak-speed magnitudes
and more consistent sequencing/timing characteristics
(\citeproc{ref-cheetham2008}{Cheetham et al. 2008}). Zheng and
colleagues' three-dimensional comparison of professional and amateur
swings found systematic temporal and spatial differences, with amateurs
deviating progressively more from professional norms as handicap
increased (\citeproc{ref-zheng2008}{Zheng et al. 2008}). Meister and
colleagues provide elite rotational benchmarks (pelvis-thorax separation
and its timing) against which amateurs can be compared
(\citeproc{ref-meister2011}{Meister et al. 2011}), and Lindsay, Mantrop
and Vandervoort review the broader catalogue of skill-level
biomechanical differences (\citeproc{ref-lindsay2008}{Lindsay, Mantrop,
and Vandervoort 2008}).

Three refinements prevent an overly tidy reading:

\begin{enumerate}
\def\labelenumi{\arabic{enumi}.}
\item
  \textbf{Sequence order is common but not universal --- even among the
  best.} Skilled golfers whose peak ordering deviates from the canonical
  pelvis-thorax-arm-club pattern are documented, and recent work
  examines how alternative sequence patterns among elite players relate
  to clubhead speed (\citeproc{ref-lee2026}{Lee, Ehrlich, and Cain
  2026}). On-plane analyses of the ``axle-chain'' system in 66 skilled
  golfers found clubhead-speed associations concentrated in specific
  angular-motion characteristics rather than in gross ordering
  (\citeproc{ref-madrid2020}{Madrid et al. 2020}).
\item
  \textbf{Within-player timing differences do not straightforwardly
  explain shot quality.} Comparing well-timed and mistimed shots
  \emph{within} highly skilled players, Neal and colleagues found no
  significant differences in the timing lags between peak angular speeds
  of contiguous segments (\citeproc{ref-neal2007}{Neal et al. 2007}) ---
  the felt sense of ``timing'' lives somewhere other than the inter-peak
  lags of a five-segment model.
\item
  \textbf{Method sensitivity.} As noted in Section~\ref{sec-energyflow},
  the identified sequence depends on component-selection methodology
  (\citeproc{ref-marsan2019}{Marsan et al. 2019}), so cross-study
  comparisons of ordering statistics carry real uncertainty.
\end{enumerate}

What survives these refinements: \emph{skilled golfers generate larger
proximal speeds, separate transition from release more distinctly, and
transmit impulse/energy to the club more efficiently; merely exhibiting
a P→D peak order is neither unique to skill nor sufficient for it}
(\citeproc{ref-cheetham2008}{Cheetham et al. 2008};
\citeproc{ref-zheng2008}{Zheng et al. 2008};
\citeproc{ref-frontiers2026footground}{{``Foot--Ground Interaction and
Clubhead Speed: Impulse-Based Energy Transfer as the Key Mechanism in
the Golf Swing''} 2026}; \citeproc{ref-lee2026}{Lee, Ehrlich, and Cain
2026}).

\subsection{Where the transfer occurs}\label{where-the-transfer-occurs}

Synthesizing Section~\ref{sec-workbudget} and Section~\ref{sec-skill}
with Section~\ref{sec-mechanics}: in a skilled swing, the dominant
energy \emph{generation} occurs in the trunk-hip complex during the
first two-thirds of the downswing (\citeproc{ref-nesbit2005work}{Nesbit
and Serrano 2005}); the dominant \emph{transfer to the club} occurs in
the final phase, through the wrist interface, by a mechanism that is
substantially parametric/geometric (wrist-cock release under centrifugal
loading, hub-path tightening) with a modest late active wrist
contribution (\citeproc{ref-white2006}{White 2006};
\citeproc{ref-miura2001}{Miura 2001};
\citeproc{ref-sprigings2000}{Sprigings and Neal 2000};
\citeproc{ref-nesbit2009hub}{Nesbit and McGinnis 2009}). Less-skilled
players truncate exactly this separation: early release (``casting'')
moves energy into the club before proximal speed has peaked, producing
the lower distal peak speeds and transfer ratios observed in amateur
cohorts (\citeproc{ref-cheetham2008}{Cheetham et al. 2008};
\citeproc{ref-zheng2008}{Zheng et al. 2008};
\citeproc{ref-milburn1982}{Milburn 1982}).

\section{The timing question, stated precisely}\label{sec-timing}

\subsection{Two strategies}\label{two-strategies}

Let \(E_p(t)\) be the kinetic energy of the proximal subsystem (trunk
and arms, with the club riding folded) over a downswing
\(t \in [t_0, t_{\mathrm{imp}}]\), and define the transfer functional
\(P_{\mathrm{w}}(t)\) as the inter-segment power crossing the wrist
interface (joint-force power plus moment power, in the sense of
Section~\ref{sec-rw}). The two strategies are:

\begin{itemize}
\tightlist
\item
  \textbf{S1 (transfer early):} make \(P_{\mathrm{w}}(t)\) large as soon
  after transition as possible; the club begins acquiring its final
  energy early.
\item
  \textbf{S2 (retain early, release late):} keep \(P_{\mathrm{w}}(t)\)
  small early (the wrist transmits centripetal force but little power)
  while proximal torques maximize \(E_p\); then allow/drive a large
  \(P_{\mathrm{w}}\) concentrated in the final phase.
\end{itemize}

The double-pendulum, delayed-release and optimal-control literatures
reviewed in Section~\ref{sec-mechanics} uniformly favor S2 within their
model classes (\citeproc{ref-milburn1982}{Milburn 1982};
\citeproc{ref-pickering1999}{Pickering and Vickers 1999};
\citeproc{ref-sprigings2000}{Sprigings and Neal 2000};
\citeproc{ref-sprigings2002}{Sprigings and MacKenzie 2002};
\citeproc{ref-sharp2009}{Sharp 2009}; \citeproc{ref-white2006}{White
2006}). The mechanistic reasons compose into a coherent account:

\begin{enumerate}
\def\labelenumi{\arabic{enumi}.}
\tightlist
\item
  \textbf{Inertia management.} With the wrist cocked, the system's
  moment of inertia about the hub is small, so proximal torque buys
  angular velocity cheaply; early release raises the inertia while
  proximal torque is still accelerating the system
  (\citeproc{ref-white2006}{White 2006}; \citeproc{ref-jorgensen1999}{T.
  P. Jorgensen 1999}).
\item
  \textbf{Actuator physiology.} Torque-velocity limits mean the wrist's
  small muscles add meaningful power only while wrist angular velocity
  is modest; early wrist torque squanders that window and can even
  require \emph{negative} (restraining) torque later
  (\citeproc{ref-sprigings2000}{Sprigings and Neal 2000}).
\item
  \textbf{Parametric transfer needs developed centrifugal loading.} The
  passive outward energy flow at release, and the inward-pull gain,
  scale with the centrifugal field established by proximal rotation ---
  weak early, strong late (\citeproc{ref-white2006}{White 2006};
  \citeproc{ref-miura2001}{Miura 2001}).
\item
  \textbf{Interaction torques do the late work for free.}
  Motion-dependent torques accelerate the distal segment and decelerate
  the proximal one automatically once release begins; the golfer's job
  is to \emph{time} this cascade, not to muscle it
  (\citeproc{ref-putnam1993}{Putnam 1993}).
\end{enumerate}

Empirically, S2-like signatures separate skill groups
(Section~\ref{sec-skill}), while the within-expert null result on
inter-peak lags (\citeproc{ref-neal2007}{Neal et al. 2007}) suggests
that once S2 is grossly achieved, residual performance variance lives in
finer variables (release profile shape, hub-path geometry,
impulse-transfer efficiency) --- exactly what the counterfactual
framework below is designed to expose.

\subsection{What would falsify S2?}\label{what-would-falsify-s2}

To keep the question scientific rather than rhetorical, we state what
evidence would count against the retain-early/release-late hypothesis:

\begin{itemize}
\tightlist
\item
  an optimal-control solution, in a validated 3D model with
  physiological actuator limits, whose optimal wrist-interface power is
  front-loaded;
\item
  measured expert swings in which early wrist-interface power
  (normalized) is \emph{positively} associated with clubhead speed
  across players;
\item
  counterfactual (ZTCF/ZVCF) analyses showing that in faster swings the
  control term dominates club acceleration \emph{early} rather than
  late.
\end{itemize}

The proposed experiments of Section~\ref{sec-proposed-experiments}
operationalize these tests.

\section{Counterfactual decompositions: ZTCF and ZVCF}\label{sec-ztcf}

\subsection{Definitions}\label{definitions}

Write the chain's equation of motion in standard manipulator form:

\begin{equation}\phantomsection\label{eq-eom}{
M(q)\,\ddot q \;=\; \tau \;-\; b(q, \dot q),
\qquad
b(q,\dot q) \;=\; C(q,\dot q)\,\dot q \;+\; g(q) \;+\; d(\dot q),
}\end{equation}

with inertia matrix \(M\), applied generalized torques \(\tau\),
Coriolis and centripetal terms \(C\dot q\), gravity \(g\), and
velocity-dependent dissipation \(d\). Two counterfactual accelerations
are defined \emph{pointwise} at each measured state
\((q, \dot q, \tau)\) of a swing:

\begin{equation}\phantomsection\label{eq-ztcf}{
\ddot q_{\mathrm{ZTCF}}
  \;=\; M(q)^{-1}\,\bigl(-\,b(q,\dot q)\bigr),
}\end{equation}

\begin{equation}\phantomsection\label{eq-zvcf}{
\ddot q_{\mathrm{ZVCF}}
  \;=\; M(q)^{-1}\,\bigl(\tau \;-\; b(q,\mathbf{0})\bigr)
  \;=\; M(q)^{-1}\,\bigl(\tau \;-\; g(q)\bigr),
}\end{equation}

where the simplification in Equation~\ref{eq-zvcf} holds because the
Coriolis term (quadratic in velocity) and viscous dissipation (linear in
velocity) vanish at \(\dot q = \mathbf 0\). The ZTCF
(Equation~\ref{eq-ztcf}) answers \emph{``what would the system do right
now if every actuator switched off?''} --- it is the \textbf{drift}
field of the affine control system, containing gravity plus all
motion-dependent interaction dynamics. The ZVCF (Equation~\ref{eq-zvcf})
answers \emph{``what would the currently applied torques, together with
gravity, produce from a standstill in this pose?''} --- it removes every
motion-dependent term but is \emph{not} the pure control contribution,
since gravity remains in Equation~\ref{eq-zvcf}; the pure torque term is
\(M^{-1}\tau\) in the split below. Because Equation~\ref{eq-eom} is
affine in \(\tau\), the actual acceleration superposes exactly:

\begin{equation}\phantomsection\label{eq-superposition}{
\ddot q \;=\;
\underbrace{M^{-1}(-b(q,\dot q))}_{\text{drift}\;=\;\text{ZTCF}}
\;+\;
\underbrace{M^{-1}\tau}_{\text{control}},
}\end{equation}

an identity enforced to numerical tolerance by the repository's test
suite (Section~\ref{sec-repo-validation}).

\subsection{Two operational variants}\label{two-operational-variants}

The repository implements two distinct operationalizations, and the
difference matters for interpretation:

\textbf{Pointwise (instantaneous) counterfactuals.} The engine-agnostic
Python implementation
(\href{https://github.com/D-sorganization/UpstreamDrift/blob/main/src/shared/python/simulation_backends/ztcf_zvcf.py}{\texttt{simulation\_backends/ztcf\_zvcf.py}})
evaluates Equation~\ref{eq-ztcf} and Equation~\ref{eq-zvcf} at each
sample of an already-measured trajectory --- explicitly \emph{not} a
forward-integrated rollout. Each sample answers a local ``what if'';
this is the right tool for drift/control \emph{time profiles} along a
real swing.

\textbf{Re-simulation counterfactuals.} The original MATLAB/Simscape
pipeline
(\href{https://github.com/D-sorganization/UpstreamDrift/blob/main/src/engines/Simscape_Multibody_Models/2D_Golf_Model/matlab_optimized/core/simulation/run_ztcf_simulation.m}{\texttt{run\_ztcf\_simulation.m}})
re-simulates the full model with a torque ``killswitch'' applied at time
\(t\), for a sweep of \(t\) values, tabulating the counterfactual state
reached with torques off. The pipeline forms the difference table
\(\mathrm{DELTA} = \mathrm{BASE} - \mathrm{ZTCF}\)
(\href{https://github.com/D-sorganization/UpstreamDrift/blob/main/src/engines/Simscape_Multibody_Models/2D_Golf_Model/matlab_optimized/core/processing/process_data_tables.m}{\texttt{process\_data\_tables.m}}),
attributing to active control everything the passive system would not
have done. A companion ZVCF model applies the measured torques to the
same pose with velocities zeroed
(\href{https://github.com/D-sorganization/UpstreamDrift/blob/main/src/engines/Simscape_Multibody_Models/2D_Golf_Model/matlab/Model\%20Output/Scripts/SCRIPT_ZVCF_GENERATOR.m}{\texttt{SCRIPT\_ZVCF\_GENERATOR.m}});
the ZVCF model is run without gravity so that, when its samples are set
against the ZTCF (which retains gravity), the gravitational contribution
is not double-counted. Note that this decomposition is meaningful
\emph{pointwise}, at matched states, in the sense of
Equation~\ref{eq-superposition}: the ZVCF script records one sample per
zero-velocity pose, and separately integrated counterfactual rollouts do
not superpose at the trajectory level in a nonlinear multibody system,
because \(M\) and \(b\) evolve differently once the counterfactual state
diverges. The MATLAB pipeline additionally computes, for every joint and
every table (BASE, ZTCF, DELTA), the linear work
\(\int \mathbf F\cdot\mathbf v\,dt\), angular work
\(\int \boldsymbol\tau\cdot\boldsymbol\omega\,dt\), total work and
power, and the fractional contributions ZTCF/BASE and DELTA/BASE
(\href{https://github.com/D-sorganization/UpstreamDrift/blob/main/src/engines/Simscape_Multibody_Models/2D_Golf_Model/matlab_optimized/core/processing/calculate_work_impulse.m}{\texttt{calculate\_work\_impulse.m}},
\href{https://github.com/D-sorganization/UpstreamDrift/blob/main/src/engines/Simscape_Multibody_Models/2D_Golf_Model/matlab_optimized/core/processing/calculate_total_work_power.m}{\texttt{calculate\_total\_work\_power.m}})
--- precisely the Robertson-Winter-style joint-interface accounting
(Section~\ref{sec-rw}) that the timing question needs, including a
kinematic-sequence plot computed \emph{separately for the BASE, ZTCF and
DELTA components}
(\href{https://github.com/D-sorganization/UpstreamDrift/blob/main/src/engines/Simscape_Multibody_Models/2D_Golf_Model/matlab/Scripts/_ZTCF\%20Scripts/SCRIPT_324_PLOT_ZTCF_KinematicSequence.m}{\texttt{SCRIPT\_324\_PLOT\_ZTCF\_KinematicSequence.m}}).

\subsection{The counterfactual form of the timing
hypothesis}\label{the-counterfactual-form-of-the-timing-hypothesis}

The ZTCF/ZVCF language makes the S2 hypothesis crisp. Along a measured
downswing, project the drift and control accelerations onto the club's
speed-producing direction and integrate the associated powers. S2
predicts, for high-performing swings:

\begin{itemize}
\tightlist
\item
  \textbf{Early downswing:} the \emph{control} term should dominate
  energy input at proximal joints, while at the wrist the control
  contribution to club angular acceleration should be small or negative
  (restraining lag against the developing centrifugal field). The club's
  ZTCF profile should show the drift field already ``wants'' to open the
  wrist --- control is \emph{holding energy back}.
\item
  \textbf{Late downswing:} the \emph{drift} term (parametric transfer,
  centrifugal release) should dominate club acceleration, with a brief
  late window of positive wrist control power
  (\citeproc{ref-sprigings2000}{Sprigings and Neal 2000}). The DELTA
  (active) share of club-side work should concentrate before release;
  the ZTCF (passive) share should account for most club energy gain
  after release (\citeproc{ref-white2006}{White 2006}).
\end{itemize}

This pattern is stated here as a \emph{hypothesis to be quantified}, not
a result: the substantiating figures and tables are the deliverable of
issue
\href{https://github.com/D-sorganization/UpstreamDrift/issues/8377}{\#8377}.
Conversely, if high-performing swings show early-dominant wrist control
power, S2 is falsified in exactly the terms of Section~\ref{sec-timing}.

\section{What UpstreamDrift computes today}\label{sec-repo}

This section inventories the repository's relevant capabilities; links
point to the \texttt{main} branch.

\subsection{Models}\label{models}

\begin{longtable}[]{@{}
  >{\raggedright\arraybackslash}p{(\linewidth - 4\tabcolsep) * \real{0.3333}}
  >{\raggedright\arraybackslash}p{(\linewidth - 4\tabcolsep) * \real{0.3333}}
  >{\raggedright\arraybackslash}p{(\linewidth - 4\tabcolsep) * \real{0.3333}}@{}}
\caption{Golf swing models available in the
repository.}\label{tbl-models}\tabularnewline
\toprule\noalign{}
\begin{minipage}[b]{\linewidth}\raggedright
Model
\end{minipage} & \begin{minipage}[b]{\linewidth}\raggedright
Location
\end{minipage} & \begin{minipage}[b]{\linewidth}\raggedright
Notes
\end{minipage} \\
\midrule\noalign{}
\endfirsthead
\toprule\noalign{}
\begin{minipage}[b]{\linewidth}\raggedright
Model
\end{minipage} & \begin{minipage}[b]{\linewidth}\raggedright
Location
\end{minipage} & \begin{minipage}[b]{\linewidth}\raggedright
Notes
\end{minipage} \\
\midrule\noalign{}
\endhead
\bottomrule\noalign{}
\endlastfoot
Double pendulum (arms + club, clubhead tip mass) &
\href{https://github.com/D-sorganization/UpstreamDrift/blob/main/src/shared/python/pendulum_simulator/physics.py}{\texttt{pendulum\_simulator/physics.py}}
& Lagrangian dynamics; viscous + Coulomb friction; analytic \(M\),
\(C\), \(g\) \\
Triple pendulum &
\href{https://github.com/D-sorganization/UpstreamDrift/blob/main/src/shared/python/pendulum_simulator/physics_triple.py}{\texttt{pendulum\_simulator/physics\_triple.py}}
& Three-segment planar chain \\
8-DOF closed-loop golfer (two arms + shared club) &
\href{https://github.com/D-sorganization/UpstreamDrift/blob/main/src/shared/python/pendulum_simulator/physics_golfer.py}{\texttt{pendulum\_simulator/physics\_golfer.py}}
& Holonomic loop constraints, constraint-consistent dynamics; JAX
variant \\
Backend-neutral golf model (ODE / MuJoCo / MJX / Warp) &
\href{https://github.com/D-sorganization/UpstreamDrift/blob/main/src/shared/python/simulation_backends/README.md}{\texttt{simulation\_backends/}}
& One parameter set rendered to every backend; shared trace schema \\
MuJoCo humanoid golfer &
\href{https://github.com/D-sorganization/UpstreamDrift/tree/main/src/engines/physics_engines/mujoco/python/mujoco_humanoid_golf}{\texttt{mujoco\_humanoid\_golf/}}
& Full-body model with club \\
Simscape 2D/3D golf models &
\href{https://github.com/D-sorganization/UpstreamDrift/tree/main/src/engines/Simscape_Multibody_Models}{\texttt{Simscape\_Multibody\_Models/}}
& Source of the original ZTCF/ZVCF methodology \\
Flexible shaft (FEM) &
\href{https://github.com/D-sorganization/UpstreamDrift/blob/main/src/shared/python/physics/flexible_shaft.py}{\texttt{physics/flexible\_shaft.py}}
& Beam model; strain energy not yet in the energy budget \\
\end{longtable}

\subsection{Counterfactual stack}\label{counterfactual-stack}

The normative contract lives in the engine-core dynamics interface
(\href{https://github.com/D-sorganization/UpstreamDrift/blob/main/src/shared/python/engine_core/_dynamics_interface.py}{\texttt{\_dynamics\_interface.py}}):
\texttt{compute\_drift\_acceleration},
\texttt{compute\_control\_acceleration}, \texttt{compute\_ztcf},
\texttt{compute\_zvcf}, with the superposition identity
(Equation~\ref{eq-superposition}) stated as a contract. Engine
implementations exist for MuJoCo, Drake, Pinocchio, OpenSim, MyoSuite,
the analytic pendulum engines, and the Simscape bridge; the capability
matrix is documented in
\href{https://github.com/D-sorganization/UpstreamDrift/blob/main/docs/engines/engine_capabilities.md}{\texttt{engine\_capabilities.md}}.
On top of these sit:

\begin{itemize}
\tightlist
\item
  the pointwise engine-agnostic evaluator with HDF5 persistence
  (\href{https://github.com/D-sorganization/UpstreamDrift/blob/main/src/shared/python/simulation_backends/ztcf_zvcf.py}{\texttt{ztcf\_zvcf.py}});
\item
  a MuJoCo delta-based counterfactual analyzer reporting
  torque-attributed and velocity-attributed acceleration effects
  (\href{https://github.com/D-sorganization/UpstreamDrift/blob/main/src/engines/physics_engines/mujoco/python/mujoco_humanoid_golf/counterfactuals.py}{\texttt{counterfactuals.py}});
\item
  a Newton-Euler ZTCF force module with an analytic single-pendulum
  validation identity
  (\href{https://github.com/D-sorganization/UpstreamDrift/blob/main/src/shared/python/biomechanics/ztcf.py}{\texttt{biomechanics/ztcf.py}});
\item
  induced-acceleration analyzers (gravity / drift / control / total
  components)
  (\href{https://github.com/D-sorganization/UpstreamDrift/blob/main/src/engines/physics_engines/mujoco/python/mujoco_humanoid_golf/rigid_body_dynamics/induced_acceleration.py}{\texttt{induced\_acceleration.py}});
\item
  an orchestration layer with capability gating and a REST/UI surface
  (\href{https://github.com/D-sorganization/UpstreamDrift/blob/main/src/shared/python/analysis/orchestrator.py}{\texttt{analysis/orchestrator.py}});
\item
  the MATLAB re-simulation pipeline described in Section~\ref{sec-ztcf}.
\end{itemize}

\subsection{Energy and power
computations}\label{energy-and-power-computations}

\begin{itemize}
\tightlist
\item
  \textbf{Inter-segment power flow (MuJoCo):}
  \href{https://github.com/D-sorganization/UpstreamDrift/blob/main/src/engines/physics_engines/mujoco/python/mujoco_humanoid_golf/power_flow.py}{\texttt{power\_flow.py}}
  computes per-DOF joint powers \(\tau\omega\), work decomposed into
  drift and control shares, per-body kinetic and potential energies via
  center-of-mass Jacobians, per-body parent/child transfer bookkeeping,
  and an energy-conservation residual derived from the equations of
  motion.
\item
  \textbf{Joint work/power/impulse metrics (engine-agnostic):}
  \href{https://github.com/D-sorganization/UpstreamDrift/blob/main/src/shared/python/analysis/power_work_metrics.py}{\texttt{power\_work\_metrics.py}}
  (positive/negative/net work, generation-absorption splits, impulse,
  quasi-stiffness);
  \href{https://github.com/D-sorganization/UpstreamDrift/blob/main/src/shared/python/analysis/energy_metrics.py}{\texttt{energy\_metrics.py}}
  (system energy summaries).
\item
  \textbf{Conservation monitoring:}
  \href{https://github.com/D-sorganization/UpstreamDrift/blob/main/src/shared/python/physics/energy_monitor.py}{\texttt{energy\_monitor.py}}
  enforces drift bounds on total energy during integration (note: its
  potential-energy term is a documented first-order approximation,
  adequate for drift detection only).
\item
  \textbf{MATLAB energetics:} the per-joint linear/angular work and
  power decomposition across BASE/ZTCF/DELTA tables
  (Section~\ref{sec-ztcf}) --- currently the richest joint-interface
  energy accounting in the repository.
\end{itemize}

\subsection{Sequencing analytics}\label{sequencing-analytics}

\begin{itemize}
\tightlist
\item
  \href{https://github.com/D-sorganization/UpstreamDrift/blob/main/src/shared/python/biomechanics/kinematic_sequence.py}{\texttt{kinematic\_sequence.py}}:
  peak detection, ordering consistency against a caller-supplied
  expected order, timing gaps, distal/proximal speed-gain ratios,
  post-peak deceleration rates. (Deliberately method-neutral; the
  expected order is user-defined.)
\item
  \href{https://github.com/D-sorganization/UpstreamDrift/blob/main/src/shared/python/analysis/coordination_metrics.py}{\texttt{coordination\_metrics.py}}:
  vector-coding coupling angles, continuous relative phase, pairwise lag
  matrices --- quantitative sequencing measures less brittle than
  peak-ordering.
\item
  Plot renderers for sequence traces, peak-timing Gantt charts, and lag
  matrices
  (\href{https://github.com/D-sorganization/UpstreamDrift/blob/main/src/shared/python/plotting/renderers/_coordination_sequence.py}{\texttt{\_coordination\_sequence.py}}).
\item
  An approved internal research review of the kinematic-sequence metric,
  including literature benchmarks and uncertainty-reporting guidance
  (\href{https://github.com/D-sorganization/UpstreamDrift/blob/main/docs/assessments/research_reviews/kinematic_sequence.md}{\texttt{kinematic\_sequence.md}}).
\end{itemize}

\subsection{Validation evidence}\label{sec-repo-validation}

The counterfactual stack is contract-tested: the pointwise evaluator
reproduces the analytic double-pendulum ground truth to \(10^{-9}\) via
the ODE and analytical providers, and agrees with the independent MuJoCo
backend to \(10^{-7}\) when that engine is available
(\href{https://github.com/D-sorganization/UpstreamDrift/blob/main/tests/unit/simulation_backends/test_ztcf_zvcf.py}{\texttt{test\_ztcf\_zvcf.py}});
drift-plus-control reconstructs the measured acceleration
(\href{https://github.com/D-sorganization/UpstreamDrift/blob/main/tests/acceptance/test_counterfactual_experiments.py}{\texttt{test\_counterfactual\_experiments.py}});
work-energy-theorem and conservation-law property tests run in CI
(\href{https://github.com/D-sorganization/UpstreamDrift/tree/main/tests/integration/conservation_laws}{\texttt{tests/integration/conservation\_laws/}});
and a cross-engine comparison service produces divergence-annotated
ZTCF/ZVCF panels between backends
(\href{https://github.com/D-sorganization/UpstreamDrift/blob/main/docs/simulation_backends/cross_engine_comparison.md}{\texttt{cross\_engine\_comparison.md}}).
The project's design guidelines make the drift-control decomposition and
ZTCF/ZVCF availability normative across engines
(\href{https://github.com/D-sorganization/UpstreamDrift/blob/main/docs/project_design_guidelines.qmd}{\texttt{project\_design\_guidelines.qmd}}).

\subsection{Gaps relevant to this article}\label{sec-repo-gaps}

\begin{longtable}[]{@{}
  >{\raggedright\arraybackslash}p{(\linewidth - 2\tabcolsep) * \real{0.5000}}
  >{\raggedright\arraybackslash}p{(\linewidth - 2\tabcolsep) * \real{0.5000}}@{}}
\caption{Known gaps (tracked in issue
\href{https://github.com/D-sorganization/UpstreamDrift/issues/8378}{\#8378}).}\label{tbl-gaps}\tabularnewline
\toprule\noalign{}
\begin{minipage}[b]{\linewidth}\raggedright
Gap
\end{minipage} & \begin{minipage}[b]{\linewidth}\raggedright
Consequence for the timing question
\end{minipage} \\
\midrule\noalign{}
\endfirsthead
\toprule\noalign{}
\begin{minipage}[b]{\linewidth}\raggedright
Gap
\end{minipage} & \begin{minipage}[b]{\linewidth}\raggedright
Consequence for the timing question
\end{minipage} \\
\midrule\noalign{}
\endhead
\bottomrule\noalign{}
\endlastfoot
Two ZTCF operationalizations (pointwise vs re-simulation) coexist; no
Python re-simulation ZTCF & Analyses must state which variant they use;
parity between variants is untested \\
Inter-segment power transfer implemented only for MuJoCo, as joint power
\(\tau\omega\) & Wrist-interface \emph{force} power
\(\mathbf F \cdot \mathbf v\) --- central to S1/S2 --- is not yet
computed in Python \\
No Robertson-Winter segmental energy-flow accounting in Python & The
transfer functional \(P_{\mathrm w}(t)\) must currently come from the
MATLAB tables \\
\texttt{ConservationMonitor} potential energy is first-order approximate
& Do not report absolute PE from it; use engine energy accessors \\
Flexible-shaft strain energy absent from the energy budget &
Late-downswing shaft recoil energy is invisible to current accounting \\
Proximal-braking-efficiency metric is an open TODO & Deceleration-side
claims must rely on generic post-peak deceleration rates \\
\end{longtable}

\section{Synthesis: early or late?}\label{sec-synthesis}

\textbf{What is established.} (i) Skilled swings exhibit P→D ordering
with successive speed amplification across shot types and populations
(\citeproc{ref-tinmark2010}{Tinmark et al. 2010};
\citeproc{ref-cheetham2008}{Cheetham et al. 2008}). (ii) The trunk-hip
complex does most of the work; the club acquires most of its energy
late, through the wrist interface (\citeproc{ref-nesbit2005work}{Nesbit
and Serrano 2005}; \citeproc{ref-nesbit2005kinematic}{Nesbit 2005}).
(iii) In every published double-pendulum, delayed-release, and
optimal-control analysis, delaying the distal handoff increases clubhead
speed, and productive distal torque is late
(\citeproc{ref-milburn1982}{Milburn 1982};
\citeproc{ref-pickering1999}{Pickering and Vickers 1999};
\citeproc{ref-sprigings2000}{Sprigings and Neal 2000};
\citeproc{ref-sprigings2002}{Sprigings and MacKenzie 2002};
\citeproc{ref-sharp2009}{Sharp 2009}). (iv) The late transfer is
substantially passive --- parametric/geometric --- once proximal
rotation has established the centrifugal field
(\citeproc{ref-white2006}{White 2006}; \citeproc{ref-miura2001}{Miura
2001}). (v) Skill differences appear in magnitudes, transition timing,
and transfer efficiency more than in gross order
(\citeproc{ref-cheetham2008}{Cheetham et al. 2008};
\citeproc{ref-zheng2008}{Zheng et al. 2008};
\citeproc{ref-frontiers2026footground}{{``Foot--Ground Interaction and
Clubhead Speed: Impulse-Based Energy Transfer as the Key Mechanism in
the Golf Swing''} 2026}).

\textbf{What this implies for the question.} The weight of evidence
supports \textbf{S2: retain energy proximally early in the downswing and
accelerate the transfer late.} The phrase ``retain energy proximally''
deserves precision: it does not mean suppressing proximal motion ---
proximal segments should be accelerated as hard as the movement budget
allows --- it means \emph{deferring the handoff}: keeping the wrist
interface transmitting force but little power (club folded, riding with
the arms) until late, so that the energy pool is large and the
centrifugal field strong when the parametric transfer is allowed to run.
In counterfactual terms: early control effort should appear at proximal
joints and, at the wrist, partly as \emph{restraint} of the drift
field's tendency to release; late club energy gain should be
drift-dominated with a brief late window of positive wrist control power
(Section~\ref{sec-ztcf}).

\textbf{What remains genuinely open.} (a) The published S2 evidence is
model-based or observational; a \emph{direct} energetic accounting of
the wrist interface across skill levels (the \(P_{\mathrm w}(t)\) time
course) has not been reported in the literature reviewed here --- the
closest is the impulse-transfer-efficiency result of
(\citeproc{ref-frontiers2026footground}{{``Foot--Ground Interaction and
Clubhead Speed: Impulse-Based Energy Transfer as the Key Mechanism in
the Golf Swing''} 2026}), which is consistent with, but does not
resolve, the timing profile. (b) Optimal timing is plausibly equipment-
and anthropometry-dependent, and model sensitivity is untested across
engines. (c) Whether \emph{actively} restraining release (negative wrist
torque early) outperforms merely \emph{not driving} it is unresolved in
3D models with realistic actuators --- 2D analyses suggest passive
holding suffices (\citeproc{ref-jorgensen1999}{T. P. Jorgensen 1999};
\citeproc{ref-white2006}{White 2006}), while actuator-based models leave
room for an active-restraint benefit
(\citeproc{ref-sprigings2000}{Sprigings and Neal 2000}). These are the
questions the repository's toolchain can answer.

\section{Proposed experiments}\label{sec-proposed-experiments}

Each design below has a pre-registered, falsifiable outcome measure. All
are tracked in issue
\href{https://github.com/D-sorganization/UpstreamDrift/issues/8377}{\#8377};
no results exist yet, and none are reported here.

\textbf{E1 --- Torque-timing sweep (double pendulum and 8-DOF golfer).}
Drive the models of Table~\ref{tbl-models} with parameterized torque
programs: proximal torque fixed; wrist-torque onset time
\(t_{\mathrm{on}}\), sign profile (restrain-then-drive vs drive-only),
and magnitude swept on a grid. Outcome: clubhead speed at impact versus
\(t_{\mathrm{on}}\) and profile. S2 predicts a late-interior optimum of
\(t_{\mathrm{on}}\) and a non-inferior restrain-then-drive profile. The
MJX/Warp batched backends make the sweep cheap; the ODE backend is the
reference.

\textbf{E2 --- Counterfactual energy budget along swept trajectories.}
For each E1 trajectory, evaluate the pointwise ZTCF/ZVCF and
drift/control split (Equation~\ref{eq-superposition}) and integrate
drift-side and control-side club power. Outcome: the drift-vs-control
share of final club kinetic energy in early vs late downswing halves, as
a function of performance. S2 predicts the drift share of late club
energy gain rises with clubhead speed.

\textbf{E3 --- Re-simulation parity.} Implement the killswitch
re-simulation ZTCF in Python on the backend-neutral model; compare
against the pointwise ZTCF profiles and the MATLAB pipeline on matched
trajectories. Outcome: documented agreement/divergence between the two
ZTCF operationalizations, closing the top gap of Table~\ref{tbl-gaps}.

\textbf{E4 --- Wrist-interface power accounting.} Implement
Robertson-Winter segmental energy flow (joint force power plus moment
power per interface; issue
\href{https://github.com/D-sorganization/UpstreamDrift/issues/8378}{\#8378}),
then compute \(P_{\mathrm w}(t)\) for E1's best and worst performers.
Outcome: the direct time course of the transfer functional defined in
Section~\ref{sec-timing} --- the primary evidence either strategy needs.

\textbf{E5 --- Mocap replication.} Run the E2/E4 analyses on measured
swings ingested through the motion pipeline (C3D → tracked motion)
across skill levels. Outcome: observational test of the counterfactual
predictions of Section~\ref{sec-ztcf} against human data, using the
uncertainty-reporting conventions of the internal research review.

\section{Limitations}\label{sec-limitations}

This article synthesizes literature and specifies methodology; it
reports no new data, and the S2 hypothesis remains unquantified in the
repository's own terms until the experiments of
Section~\ref{sec-proposed-experiments} are executed. The
kinematic-sequence evidence base carries methodological sensitivity
(\citeproc{ref-marsan2019}{Marsan et al. 2019}); several classical
results rest on planar models whose generalization to 3D swings with
mobile hubs is a live question (\citeproc{ref-mcphee2022}{McPhee 2022});
and the work-budget percentages cited from
(\citeproc{ref-nesbit2005work}{Nesbit and Serrano 2005}) derive from
four subjects. One recent source
(\citeproc{ref-frontiers2026footground}{{``Foot--Ground Interaction and
Clubhead Speed: Impulse-Based Energy Transfer as the Key Mechanism in
the Golf Swing''} 2026}) is cited without its author list because the
publisher's site was unreachable through this environment's network
proxy; resolve the DOI for full details. Finally, the arguments here
concern clubhead \emph{speed}; accuracy, consistency, and injury-risk
trade-offs of late-transfer strategies (\citeproc{ref-hume2005}{Hume,
Keogh, and Reid 2005}; \citeproc{ref-lindsay2008}{Lindsay, Mantrop, and
Vandervoort 2008}) are outside scope.

\section{Conclusion}\label{sec-conclusion}

Proximal-to-distal sequencing in the golf swing is real, robust, and ---
as a coaching target --- chronically over-interpreted. The mechanical
literature, from Cochran and Stobbs through modern optimal control, does
not support ``pass the energy along as soon as possible.'' It supports a
subtler program: \emph{build the energy proximally with the club held
folded, defer the handoff, and let the transfer run --- mostly
parametrically, with a brief late assist --- when the centrifugal field
can pay for it}. Skilled players differ from amateurs along exactly this
axis: not in whether their segments peak in order, but in how much
energy the proximal system holds at transition and how late and how
efficiently it is handed to the club. The ZTCF/ZVCF framework turns this
qualitative account into measurable time profiles --- drift versus
control, early versus late --- and the experiment suite specified here
(epic
\href{https://github.com/D-sorganization/UpstreamDrift/issues/8373}{\#8373})
is designed to quantify the retain-early/release-late advantage or
falsify it. Either outcome would be a contribution.

\section*{Acknowledgments}\label{acknowledgments}
\addcontentsline{toc}{section}{Acknowledgments}

Prepared within the UpstreamDrift project for eventual publication in
the affinedrift repository. The inventory in Section~\ref{sec-repo}
reflects the codebase at the commit referenced by the pull request
introducing this document.

\section*{References}\label{references}
\addcontentsline{toc}{section}{References}

\phantomsection\label{refs}
\begin{CSLReferences}{1}{0}
\bibitem[\citeproctext]{ref-balzerson2016}
Balzerson, Daniel, Joydeep Banerjee, and John McPhee. 2016. {``A
Three-Dimensional Forward Dynamic Model of the Golf Swing Optimized for
Ball Carry Distance.''} \emph{Sports Engineering} 19: 237--50.
\url{https://doi.org/10.1007/s12283-016-0197-7}.

\bibitem[\citeproctext]{ref-bunn1972}
Bunn, John W. 1972. \emph{Scientific Principles of Coaching}. 2nd ed.
Englewood Cliffs, NJ: Prentice-Hall.

\bibitem[\citeproctext]{ref-cheetham2008}
Cheetham, Phillip J., Gregory A. Rose, Richard N. Hinrichs, Robert J.
Neal, Rob E. Mottram, Paul D. Hurrion, and Peter F. Vint. 2008.
{``Comparison of Kinematic Sequence Parameters Between Amateur and
Professional Golfers.''} In \emph{Science and Golf v: Proceedings of the
World Scientific Congress of Golf}, edited by Debbie Crews and Rafer
Lutz. Mesa, AZ: Energy in Motion.
\url{https://www.researchgate.net/publication/252236426}.

\bibitem[\citeproctext]{ref-cochran1968}
Cochran, Alastair, and John Stobbs. 1968. \emph{Search for the Perfect
Swing}. Philadelphia: J.~B. Lippincott.

\bibitem[\citeproctext]{ref-frontiers2026footground}
{``Foot--Ground Interaction and Clubhead Speed: Impulse-Based Energy
Transfer as the Key Mechanism in the Golf Swing.''} 2026.
\emph{Frontiers in Sports and Active Living} 8.
\url{https://doi.org/10.3389/fspor.2026.1790645}.

\bibitem[\citeproctext]{ref-han2019}
Han, Ki Hoon, Christopher Como, Jemin Kim, Sangwoo Lee, Jaewoong Kim,
Dae Kyoo Kim, and Young-Hoo Kwon. 2019. {``Effects of the Golfer--Ground
Interaction on Clubhead Speed in Skilled Male Golfers.''} \emph{Sports
Biomechanics} 18 (2): 115--34.
\url{https://pubmed.ncbi.nlm.nih.gov/31042142/}.

\bibitem[\citeproctext]{ref-hume2005}
Hume, Patria A., Justin Keogh, and Duncan Reid. 2005. {``The Role of
Biomechanics in Maximising Distance and Accuracy of Golf Shots.''}
\emph{Sports Medicine} 35 (5): 429--49.
\url{https://doi.org/10.2165/00007256-200535050-00005}.

\bibitem[\citeproctext]{ref-jorgensen1970}
Jorgensen, Theodore. 1970. {``On the Dynamics of the Swing of a Golf
Club.''} \emph{American Journal of Physics} 38 (5): 644--51.
\url{https://doi.org/10.1119/1.1976419}.

\bibitem[\citeproctext]{ref-jorgensen1999}
Jorgensen, Theodore P. 1999. \emph{The Physics of Golf}. 2nd ed. New
York: Springer / AIP Press.

\bibitem[\citeproctext]{ref-lampsa1975}
Lampsa, Michael A. 1975. {``Maximizing Distance of the Golf Drive: An
Optimal Control Study.''} \emph{Journal of Dynamic Systems, Measurement,
and Control} 97 (4): 362--67.
\url{https://asmedigitalcollection.asme.org/dynamicsystems/article-abstract/97/4/362/400950}.

\bibitem[\citeproctext]{ref-lee2026}
Lee, Junghoon, Justin Ehrlich, and Christopher Cain. 2026. {``The
Effects of Different Kinematic Sequences of the Elite Golf Swing on
Clubhead Speed.''} \emph{International Journal of Sports Science \&
Coaching}. \url{https://doi.org/10.1177/17479541261429479}.

\bibitem[\citeproctext]{ref-lindsay2008}
Lindsay, David M., Shannon Mantrop, and Anthony A. Vandervoort. 2008.
{``A Review of Biomechanical Differences Between Golfers of Varied Skill
Levels.''} \emph{International Journal of Sports Science \& Coaching} 3
(Suppl.~1): 187--97. \url{https://doi.org/10.1260/174795408785024117}.

\bibitem[\citeproctext]{ref-mackenzie2009}
MacKenzie, Sasho J., and Eric J. Sprigings. 2009. {``A Three-Dimensional
Forward Dynamics Model of the Golf Swing.''} \emph{Sports Engineering}
11 (4): 165--75. \url{https://doi.org/10.1007/s12283-009-0020-9}.

\bibitem[\citeproctext]{ref-madrid2020}
Madrid, Braden, Marco A. Avalos, Nicholas A. Levine, Noelle J. Tuttle,
Kevin A. Becker, and Young-Hoo Kwon. 2020. {``Association Between the
on-Plane Angular Motions of the Axle-Chain System and Clubhead Speed in
Skilled Male Golfers.''} \emph{Applied Sciences} 10 (17): 5728.
\url{https://doi.org/10.3390/app10175728}.

\bibitem[\citeproctext]{ref-marsan2019}
Marsan, Thibault, Maxime Bourgain, Patricia Thoreux, Christophe Sauret,
and Philippe Rouch. 2019. {``Biomechanical Analysis of the Golf Swing:
Methodological Effect of Angular Velocity Component on the
Identification of the Kinematic Sequence.''} \emph{Acta of
Bioengineering and Biomechanics} 21: 115--20.
\url{https://www.researchgate.net/publication/333377321}.

\bibitem[\citeproctext]{ref-marshall2000}
Marshall, Robert N., and Bruce C. Elliott. 2000. {``Long-Axis Rotation:
The Missing Link in Proximal-to-Distal Segmental Sequencing.''}
\emph{Journal of Sports Sciences} 18 (4): 247--54.
\url{https://doi.org/10.1080/026404100364983}.

\bibitem[\citeproctext]{ref-mcnally2018}
McNally, William, and John McPhee. 2018. {``Dynamic Optimization of the
Golf Swing Using a Six Degree-of-Freedom Biomechanical Model.''} In
\emph{Proceedings}, 2:243. 6.
\url{https://doi.org/10.3390/proceedings2060243}.

\bibitem[\citeproctext]{ref-mcphee2022}
McPhee, John. 2022. {``A Review of Dynamic Models and Measurements in
Golf.''} \emph{Sports Engineering} 25.
\url{https://doi.org/10.1007/s12283-022-00387-0}.

\bibitem[\citeproctext]{ref-meister2011}
Meister, David W., Amy L. Ladd, Erin E. Butler, Betty Zhao, Andrew P.
Rogers, Conrad J. Ray, and Jessica Rose. 2011. {``Rotational
Biomechanics of the Elite Golf Swing: Benchmarks for Amateurs.''}
\emph{Journal of Applied Biomechanics} 27 (3): 242--51.
\url{https://doi.org/10.1123/jab.27.3.242}.

\bibitem[\citeproctext]{ref-milburn1982}
Milburn, Peter D. 1982. {``Summation of Segmental Velocities in the Golf
Swing.''} \emph{Medicine and Science in Sports and Exercise} 14 (1):
60--64. \url{https://pubmed.ncbi.nlm.nih.gov/7070260/}.

\bibitem[\citeproctext]{ref-miura2001}
Miura, Ken. 2001. {``Parametric Acceleration --- the Effect of Inward
Pull of the Golf Club at Impact Stage.''} \emph{Sports Engineering} 4
(2): 75--86. \url{https://doi.org/10.1046/j.1460-2687.2001.00071.x}.

\bibitem[\citeproctext]{ref-neal2007}
Neal, Robert J., Ryan Lumsden, Mark Holland, and Bruce Mason. 2007.
{``Body Segment Sequencing and Timing in Golf.''} \emph{International
Journal of Sports Science \& Coaching} 2 (Suppl.~1).
\url{https://doi.org/10.1260/174795407789705497}.

\bibitem[\citeproctext]{ref-nesbit2005kinematic}
Nesbit, Steven M. 2005. {``A Three Dimensional Kinematic and Kinetic
Study of the Golf Swing.''} \emph{Journal of Sports Science and
Medicine} 4 (4): 499--519.
\url{https://pmc.ncbi.nlm.nih.gov/articles/PMC3899667/}.

\bibitem[\citeproctext]{ref-nesbit2009hub}
Nesbit, Steven M., and Ryan McGinnis. 2009. {``Kinematic Analyses of the
Golf Swing Hub Path and Its Role in Golfer/Club Kinetic Transfers.''}
\emph{Journal of Sports Science and Medicine} 8 (2): 235--46.
\url{https://www.jssm.org/volume08/iss2/cap/jssm-08-235.pdf}.

\bibitem[\citeproctext]{ref-nesbit2014hub}
Nesbit, Steven M., and Ryan S. McGinnis. 2014. {``Kinetic Constrained
Optimization of the Golf Swing Hub Path.''} \emph{Journal of Sports
Science and Medicine}. \url{https://pubmed.ncbi.nlm.nih.gov/25435779/}.

\bibitem[\citeproctext]{ref-nesbit2005work}
Nesbit, Steven M., and Monika Serrano. 2005. {``Work and Power Analysis
of the Golf Swing.''} \emph{Journal of Sports Science and Medicine} 4
(4): 520--33.
\url{https://www.jssm.org/volume04/iss4/cap/jssm-04-520.pdf}.

\bibitem[\citeproctext]{ref-pickering1999}
Pickering, W. M., and G. T. Vickers. 1999. {``On the Double Pendulum
Model of the Golf Swing.''} \emph{Sports Engineering} 2 (3): 161--72.
\url{https://doi.org/10.1046/j.1460-2687.1999.00028.x}.

\bibitem[\citeproctext]{ref-putnam1991}
Putnam, Carol A. 1991. {``A Segment Interaction Analysis of
Proximal-to-Distal Sequential Segment Motion Patterns.''} \emph{Medicine
and Science in Sports and Exercise} 23 (1): 130--44.
\url{https://pubmed.ncbi.nlm.nih.gov/1997807/}.

\bibitem[\citeproctext]{ref-putnam1993}
---------. 1993. {``Sequential Motions of Body Segments in Striking and
Throwing Skills: Descriptions and Explanations.''} \emph{Journal of
Biomechanics} 26 (Suppl.~1): 125--35.
\url{https://doi.org/10.1016/0021-9290(93)90084-R}.

\bibitem[\citeproctext]{ref-robertson1980}
Robertson, D. Gordon E., and David A. Winter. 1980. {``Mechanical Energy
Generation, Absorption and Transfer Amongst Segments During Walking.''}
\emph{Journal of Biomechanics} 13 (10): 845--54.
\url{https://doi.org/10.1016/0021-9290(80)90172-4}.

\bibitem[\citeproctext]{ref-sharp2009}
Sharp, Robin S. 2009. {``On the Mechanics of the Golf Swing.''}
\emph{Proceedings of the Royal Society A: Mathematical, Physical and
Engineering Sciences} 465 (2102): 551--70.
\url{https://doi.org/10.1098/rspa.2008.0304}.

\bibitem[\citeproctext]{ref-sprigings2002}
Sprigings, Eric J., and Sasho J. MacKenzie. 2002. {``Examining the
Delayed Release in the Golf Swing Using Computer Simulation.''}
\emph{Sports Engineering} 5 (1): 23--32.
\url{http://sashomackenzie.com/publications/delayedrelease.pdf}.

\bibitem[\citeproctext]{ref-sprigings2000}
Sprigings, Eric J., and Robert J. Neal. 2000. {``An Insight into the
Importance of Wrist Torque in Driving the Golfball: A Simulation
Study.''} \emph{Journal of Applied Biomechanics} 16 (4): 356--66.
\url{https://doi.org/10.1123/jab.16.4.356}.

\bibitem[\citeproctext]{ref-tinmark2010}
Tinmark, Fredrik, John Hellström, Kjartan Halvorsen, and Alf
Thorstensson. 2010. {``Elite Golfers' Kinematic Sequence in Full-Swing
and Partial-Swing Shots.''} \emph{Sports Biomechanics} 9 (4): 236--44.
\url{https://doi.org/10.1080/14763141.2010.535842}.

\bibitem[\citeproctext]{ref-white2006}
White, Rod. 2006. {``On the Efficiency of the Golf Swing.''}
\emph{American Journal of Physics} 74 (12): 1088--94.
\url{https://doi.org/10.1119/1.2346688}.

\bibitem[\citeproctext]{ref-zajac1989}
Zajac, Felix E., and Michael E. Gordon. 1989. {``Determining Muscle's
Force and Action in Multi-Articular Movement.''} \emph{Exercise and
Sport Sciences Reviews} 17: 187--230.
\url{https://pubmed.ncbi.nlm.nih.gov/2676547/}.

\bibitem[\citeproctext]{ref-zajac2002}
Zajac, Felix E., Richard R. Neptune, and Steven A. Kautz. 2002.
{``Biomechanics and Muscle Coordination of Human Walking. {P}art~{I}:
Introduction to Concepts, Power Transfer, Dynamics and Simulations.''}
\emph{Gait \& Posture} 16 (3): 215--32.
\url{https://doi.org/10.1016/S0966-6362(02)00068-1}.

\bibitem[\citeproctext]{ref-zheng2008}
Zheng, Naiquan, Steve W. Barrentine, Glenn S. Fleisig, and James R.
Andrews. 2008. {``Kinematic Analysis of Swing in Pro and Amateur
Golfers.''} \emph{International Journal of Sports Medicine} 29 (6):
487--93. \url{https://doi.org/10.1055/s-2008-1038732}.

\end{CSLReferences}




\end{document}
